\chapter{Przechodzenie po wszystkich krawędziach: teoria Eulera}
\label{ch:tg-euler}

Pierwsze dwa rozdziały dotyczyły przede wszystkim spójności i istnienia
dróg. Teraz stawiamy inne pytanie: czy można przejść przez \emph{każdą
krawędź dokładnie raz}? Problem ma historyczne źródło w zagadnieniu
mostów królewieckich. Prowadzi do kryterium całkowicie wyrażonego
przez stopnie wierzchołków, a także do konstruktywnego algorytmu.

\section{Szlaki i cykle Eulera}
\label{sec:tg-euler-def}

W tej sekcji dopuszczamy multigrafy nieskierowane, także z pętlami.
W zapisie spaceru wskazujemy kolejno zarówno wierzchołki, jak i
krawędzie: $v_0,e_1,v_1,\ldots,e_k,v_k$, gdzie końcami $e_i$ są
$v_{i-1}$ i $v_i$. Taki zapis odróżnia krawędzie równoległe.
\emph{Nośnikiem krawędziowym} grafu o niepustym zbiorze
krawędzi nazywamy podgraf o zbiorze wierzchołków złożonym
z końców wszystkich krawędzi i o pełnym zbiorze krawędzi.
Wierzchołki izolowane nie wpływają na możliwość przejścia
przez wszystkie krawędzie. Jeżeli $E=\varnothing$, nośnik
jest grafem zerowym i stosujemy osobną konwencję
z poniższej definicji.

\begin{definicja}[Szlak Eulera i cykl Eulera]
Niech $G=(V,E,\partial)$ będzie skończonym multigrafem nieskierowanym.
\emph{Szlakiem Eulera} nazywamy szlak, w którym każda krawędź
$e\in E$ występuje dokładnie raz. Jeżeli taki szlak jest zamknięty,
nazywamy go \emph{cyklem Eulera} (mimo że może wielokrotnie
odwiedzać ten sam wierzchołek). Dla $E=\varnothing$ i $V\ne\varnothing$ przyjmujemy,
że spacer długości zero, w dowolnym wierzchołku,
jest zamkniętym szlakiem Eulera. W tym szczególnym przypadku
nie należy wymagać spójności nośnika krawędziowego.
Dla $V=\varnothing$ nie istnieje żaden spacer,
nawet długości zero, a więc nie istnieje szlak Eulera.
\end{definicja}

Termin ,,cykl Eulera'' jest tradycyjny: nie jest to cykl prosty
zdefiniowany wcześniej. Przykładowo obieg dwóch trójkątów
połączonych jednym wspólnym wierzchołkiem jest cyklem Eulera,
choć wspólny wierzchołek odwiedzamy więcej niż raz.

\begin{center}
\begin{tikzpicture}[scale=.95,line cap=round,>=Stealth]
  \node[tgBlueVertex] (a) at (-2,0) {$a$};
  \node[tgBlueVertex] (b) at (-1.1,1.2) {$b$};
  \node[tgBlueVertex] (c) at (0,0) {$c$};
  \node[tgBlueVertex] (d) at (1.1,1.2) {$d$};
  \node[tgBlueVertex] (e) at (2,0) {$e$};
  \draw[tgTree,-{Stealth[length=2mm]}] (a)--(b);
  \draw[tgTree,-{Stealth[length=2mm]}] (b)--(c);
  \draw[tgTree,-{Stealth[length=2mm]}] (c)--(a);
  \draw[tgAccent,-{Stealth[length=2mm]}] (c)--(d);
  \draw[tgAccent,-{Stealth[length=2mm]}] (d)--(e);
  \draw[tgAccent,-{Stealth[length=2mm]}] (e)--(c);
  \node[text=tgNavy,font=\footnotesize] at (0,-.65)
     {$a\to b\to c\to d\to e\to c\to a$};
\end{tikzpicture}

{\footnotesize Dwa trójkąty mają wspólny wierzchołek $c$.
Kolory pokazują dwa fragmenty jednego zamkniętego szlaku Eulera.}
\end{center}

\begin{lemat}[Parzystość stopni wewnętrznych]
\label{lem:tg-euler-parity}
Niech $v_0,e_1,v_1,\dots,e_k,v_k$ będzie szlakiem
w skończonym multigrafie nieskierowanym $G$. Niech $H$ będzie podgrafem o zbiorze wierzchołków $V(G)$
i zbiorze krawędzi $\{e_1,\ldots,e_k\}$.
Jeżeli szlak jest zamknięty, to każdy wierzchołek $H$
ma parzysty stopień. Jeżeli $v_0\ne v_k$,
to stopnie $v_0,v_k$ w tym podgrafie są nieparzyste,
a wszystkie pozostałe są parzyste.
\end{lemat}

\begin{proof}
Jeżeli $k=0$, podgraf nie ma krawędzi i wszystkie
jego stopnie wynoszą zero. Niech więc $k\ge1$.
Każde pojawienie się wierzchołka jako punktu pośredniego
szlaku ma dwie incydentne krawędzie: tę, którą weszliśmy,
i tę, którą wyszliśmy. Krawędzie te są różne, ponieważ szlak
nie powtarza krawędzi. W ten sposób można pogrupować
wszystkie końce wykorzystanych krawędzi w pary.
Jeżeli początek i koniec są tym samym wierzchołkiem,
pierwszy wyjazd parujemy z ostatnim powrotem; również pętla
liczona jest jako dwa końce tej samej krawędzi. Przy różnych
końcach zostaje dokładnie jeden niesparowany koniec
przy $v_0$ i jeden przy $v_k$. To daje odpowiednie parzystości.
\end{proof}

\begin{twierdzenie}[Euler: charakterystyka zamknięta]
\label{thm:tg-euler-circuit}
Niech $G$ będzie skończonym multigrafem nieskierowanym
z co najmniej jedną krawędzią. Wówczas istnieje cykl Eulera
w $G$ wtedy i tylko wtedy, gdy
\begin{enumerate}[label=(\roman*)]
\item wszystkie wierzchołki dodatniego stopnia leżą w jednej
składowej spójnej;
\item dla każdego $v\in V(G)$ stopień $d_G(v)$ jest parzysty.
\end{enumerate}
\end{twierdzenie}

\begin{proof}
Jeżeli istnieje cykl Eulera, jego kolejne krawędzie łączą
wszystkie wierzchołki o dodatnim stopniu, a lemat
\ref{lem:tg-euler-parity} zapewnia parzystość ich stopni.

Załóżmy odwrotnie, że (i)--(ii) zachodzą. Wybieramy wierzchołek
$v_0$ dodatniego stopnia i rozpoczynamy szlak, wybierając
za każdym razem dowolną \emph{niewykorzystaną} krawędź.
Ponieważ graf jest skończony, proces musi się zatrzymać.
Nie może zakończyć się w wierzchołku $w\ne v_0$:
przed ewentualnym zakończeniem w $w$ liczba wykorzystanych
końców krawędzi incydentnych z $w$ jest nieparzysta,
a całkowity stopień $w$ jest parzysty. Pozostaje zatem
niewykorzystana krawędź wychodząca z $w$.
Otrzymaliśmy pewien zamknięty szlak $C$.

Jeżeli $C$ zawiera wszystkie krawędzie, dowód jest zakończony.
W przeciwnym razie istnieje niewykorzystana krawędź
przylegająca do wierzchołka $C$. Istotnie, jeżeli pewna niewykorzystana krawędź ma koniec
na $C$, otrzymujemy tezę. W przeciwnym razie wybieramy
drogę od wierzchołka $C$ do końca dowolnej niewykorzystanej
krawędzi; istnieje ona dzięki spójności nośnika.
Pierwsza krawędź tej drogi opuszczająca zbiór wierzchołków
$C$ nie może należeć do $C$, ponieważ obydwa końce każdej
krawędzi $C$ leżą w tym zbiorze. Jest więc poszukiwaną
niewykorzystaną krawędzią. Niech $w$ będzie takim wierzchołkiem.
W multigrafie złożonym z \emph{niewykorzystanych} krawędzi
wszystkie stopnie nadal są parzyste: od pierwotnego
parzystego stopnia odjęliśmy parzysty stopień pochodzący
z zamkniętego szlaku $C$. Możemy więc, zaczynając w $w$,
zbudować kolejny niepusty szlak zamknięty $C'$ złożony
wyłącznie z niewykorzystanych krawędzi.

W miejscu wystąpienia $w$ na $C$ wstawiamy cały obieg $C'$
i wracamy następnie do dalszej części $C$. Powstaje dłuższy
zamknięty szlak, nadal bez powtarzania krawędzi.
Przy każdym takim kroku zużywamy co najmniej jedną nową
krawędź. Po skończonej liczbie kroków uzyskamy cykl Eulera.
\end{proof}

\begin{uwaga}[Algorytm Hierholzera]
Dowód jest algorytmem: budujemy zamknięty szlak, wyszukujemy
na nim wierzchołek incydentny z niewykorzystaną krawędzią,
a następnie wplatamy nowy zamknięty szlak. Aby uzyskać czas $O(|V|+|E|)$, trzeba unikać ponownego
przeglądania całego obiegu przy każdym wplataniu.
Przechowujemy obieg jako listę dwukierunkową wystąpień
wierzchołków i krawędzi, a wystąpienia oczekujące na
sprawdzenie w osobnej kolejce. Przy każdym wierzchołku
utrzymujemy wskaźnik na pierwszą jeszcze niesprawdzoną
pozycję listy incydencji. Pozycje wykorzystanych krawędzi
pomijamy; każdą krawędź oznaczamy jako wykorzystaną
niezależnie od tego, przez który jej koniec ją znaleziono.
Dla pobranego wystąpienia $w$ budujemy i wplatamy kolejne
zamknięte szlaki zaczynające się w $w$, dopóki pozostaje
tam niewykorzystana krawędź. Nowe wystąpienia dodajemy
do kolejki. Gdy przy $w$ nie ma już niewykorzystanych
krawędzi, nie pojawią się one później. Każde wystąpienie
jest sprawdzane najwyżej raz, a każde wplatanie jest
stałoczasowe poza samą konstrukcją nowego szlaku.
Lista incydencji ma łącznie $2|E|$ pozycji (także dla
pętli), a lista obiegu $O(|E|)$ pozycji. Po wyczerpaniu
kolejki nie może pozostać niewykorzystana krawędź,
gdyż spójność nośnika wymuszałaby jej incydencję
z pewnym wystąpieniem w obiegu. Stąd podana złożoność,
wliczając przygotowanie list i wyszukanie korzenia. Sama zasada ,,idź dowolną krawędzią''
bez etapu wplatania cykli nie gwarantuje natychmiastowego
wykorzystania wszystkich krawędzi.
\end{uwaga}

\begin{przyklad}[Wplatanie obiegów w multigrafie]
\label{ex:tg-hierholzer-multigraf}
Niech $G$ ma wierzchołki $a,b,c$ oraz pięć krawędzi:
$e_1,e_2$ o końcach $a,b$, pętlę $e_3$ przy $b$
oraz $e_4,e_5$ o końcach $b,c$. Stopnie wynoszą
$d_G(a)=2$, $d_G(b)=6$ i $d_G(c)=2$, gdyż pętla
wnosi dwa do stopnia $b$. Graf jest spójny, a wszystkie
jego stopnie są parzyste.

Rozpoczynamy w $a$ i wybieramy najpierw $e_1$, a następnie
$e_2$. Otrzymujemy zamknięty szlak $C=(a,e_1,b,e_2,a)$.
Nie można go już przedłużyć od jego końca $a$, chociaż
trzy krawędzie pozostają niewykorzystane. Wybieramy
wystąpienie $b$ w $C$. W multigrafie pozostałych krawędzi
możemy zbudować zamknięty szlak
$C'=(b,e_3,b,e_4,c,e_5,b)$.
Wplatając $C'$ w miejsce wystąpienia $b$ w $C$,
otrzymujemy
\[
(a,e_1,b,e_3,b,e_4,c,e_5,b,e_2,a).
\]
Szlak jest zamknięty, a lista jego krawędzi to
$e_1,e_3,e_4,e_5,e_2$, czyli każda z pięciu krawędzi
występuje dokładnie raz. Jest to cykl Eulera.
Przykład pokazuje, dlaczego samo przedłużanie jednego
szlaku aż do zatrzymania nie wystarcza oraz dlaczego
należy odróżniać równoległe krawędzie.
\end{przyklad}

\begin{twierdzenie}[Euler: charakterystyka otwarta]
\label{thm:tg-euler-trail}
Niech $G$ będzie skończonym multigrafem nieskierowanym
z co najmniej jedną krawędzią, którego wszystkie wierzchołki
dodatniego stopnia należą do jednej składowej spójnej.
Wówczas $G$ ma szlak Eulera o różnych końcach wtedy
i tylko wtedy, gdy dokładnie dwa wierzchołki mają
nieparzysty stopień. Końcami dowolnego takiego szlaku
są właśnie te dwa wierzchołki.
\end{twierdzenie}

\begin{proof}
Konieczność daje lemat \ref{lem:tg-euler-parity}.
Jeśli dokładnie dwa wierzchołki $s,t$ mają nieparzysty
stopień, dołączamy \emph{nową} krawędź $e_*$ o końcach $s,t$,
nawet gdy były już połączone krawędzią. Otrzymany multigraf
ma wszystkie stopnie parzyste oraz spójny nośnik krawędziowy.
Na mocy twierdzenia \ref{thm:tg-euler-circuit} ma cykl Eulera.
W razie potrzeby odwracamy kierunek obiegu tak, aby
przejście przez $e_*$ prowadziło od $t$ do $s$.
Przesuwamy cyklicznie początek zapisu obiegu do wystąpienia
$s$ bezpośrednio po $e_*$. Wtedy $e_*$ jest ostatnią
krawędzią zapisu. Usuwając ją i końcowe wystąpienie $s$,
otrzymujemy szlak od $s$ do $t$, w którym każda
oryginalna krawędź występuje dokładnie raz.
\end{proof}

\begin{przyklad}[Mosty królewieckie]
Model historycznego układu siedmiu mostów ma cztery
wierzchołki odpowiadające obszarom lądu. Ich stopnie
wynoszą $3,3,3,5$; suma $14$ odpowiada siedmiu krawędziom,
zgodnie z lematem o uściskach dłoni. Ponieważ wierzchołków
nieparzystych jest cztery, nie istnieje ani zamknięty,
ani otwarty szlak Eulera. Problem nie zależy od długości
mostów ani od ich rysunkowego położenia.
\end{przyklad}

\begin{center}
\begin{tikzpicture}[scale=.95,line cap=round]
 \node[tgBlueVertex] (a) at (-2,0) {$A$};
 \node[tgBlueVertex] (b) at (2,0) {$B$};
 \node[tgRedVertex] (c) at (0,1.45) {$C$};
 \node[tgRedVertex] (d) at (0,-1.45) {$D$};
 \draw[tgEdge] (a) to[bend left=19] (c);
 \draw[tgEdge] (a) to[bend right=19] (c);
 \draw[tgEdge] (b) to[bend left=19] (c);
 \draw[tgEdge] (b) to[bend right=19] (c);
 \draw[tgAccent] (a)--(d);
 \draw[tgAccent] (b)--(d);
 \draw[tgEdge] (c)--(d);
 \node[font=\footnotesize,text=tgNavy] at (0,-2.10)
    {$d(A)=3,\ d(B)=3,\ d(C)=5,\ d(D)=3$};
\end{tikzpicture}

{\footnotesize Jeden z abstrakcyjnych modeli problemu siedmiu mostów;
krawędzie równoległe odpowiadają różnym mostom.}
\end{center}

\section{Orientacje krawędzi i skierowane obiegi Eulera}
\label{sec:tg-orientacje}

\begin{definicja}[Graf skierowany i orientacja]
\emph{Graf skierowany} to para $(V,A)$, gdzie $V$ jest
skończonym zbiorem wierzchołków, a $A\subseteq V\times V$
jest zbiorem łuków; łuk $(u,v)$ zapisujemy jako $u\to v$.
\emph{Orientacja} grafu prostego nieskierowanego $G=(V,E)$
powstaje przez zastąpienie każdej krawędzi $\{u,v\}\in E$
dokładnie jednym z łuków $(u,v)$ lub $(v,u)$.
\emph{Multigraf skierowany} to układ $(V,A,\partial_-,\partial_+)$,
gdzie $V,A$ są skończonymi zbiorami wierzchołków i łuków,
a $\partial_-,\partial_+:A\to V$ przypisują łukowi
odpowiednio jego początek i koniec. Różne łuki mogą
mieć ten sam początek i koniec. Graf skierowany $(V,A)$
jest szczególnym przypadkiem tego modelu, z odwzorowaniami
będącymi rzutami na współrzędne pary.
Dla multigrafu nieskierowanego orientacja zachowuje odrębne
identyfikatory krawędzi i przypisuje każdej z nich
uporządkowaną parę jej końców. Pętla daje łuk od
wierzchołka do niego samego.
Stopień wejściowy $d^-(v)$ jest liczbą łuków kończących się
w $v$, a stopień wyjściowy $d^+(v)$ jest liczbą łuków
zaczynających się w $v$, z uwzględnieniem krotności.
Pętla wnosi po jednym do obu stopni.
\end{definicja}

Dla multigrafów skierowanych rozważamy spacery zgodne
z kierunkami łuków. Skierowany szlak Eulera to taki
spacer, który wykorzystuje każdy łuk dokładnie raz.
Dla zbioru pustego łuków przyjmujemy tę samą konwencję
co w przypadku nieskierowanym: szlak długości zero istnieje
dokładnie wtedy, gdy zbiór wierzchołków jest niepusty. Graf skierowany nazywamy \emph{silnie spójnym}, gdy
ma niepusty zbiór wierzchołków i każde dwa wierzchołki
można połączyć skierowanymi drogami w obie strony. Przez spójność \emph{słabą} rozumiemy spójność
grafu otrzymanego po zapomnieniu kierunków łuków.

\begin{twierdzenie}[Skierowane obiegi Eulera]
\label{thm:tg-euler-directed-circuit}
Niech $D=(V,A,\partial_-,\partial_+)$ będzie skończonym
multigrafem skierowanym
z co najmniej jednym łukiem. Wówczas ma on skierowany
zamknięty szlak Eulera wtedy i tylko wtedy, gdy
\begin{enumerate}[label=(\roman*)]
\item wszystkie wierzchołki incydentne z łukami należą do
jednej słabo spójnej składowej;
\item $d^+(v)=d^-(v)$ dla każdego $v\in V$.
\end{enumerate}
\end{twierdzenie}

\begin{proof}
W skierowanym zamkniętym szlaku każde wejście do wierzchołka
jest sparowane z wyjściem, więc otrzymujemy (ii).
Ponieważ szlak zawiera wszystkie łuki, zachodzi też (i).

Odwrotnie, wybieramy wierzchołek incydentny z łukiem.
Z warunku (ii) wynika, że ma on dodatni stopień wyjściowy.
Przedłużamy skierowany szlak bez powtarzania łuków,
dopóki istnieje niewykorzystany łuk wychodzący z jego końca.
Proces kończy się, gdyż zbiór łuków jest skończony.
Jeżeli zatrzymałby się w wierzchołku innym niż początek,
to miałby tam o jedno wykorzystane wejście więcej niż
wykorzystanych wyjść. Z równości $d^+=d^-$ wynikałoby
istnienie niewykorzystanego wyjścia, sprzeczność.
Otrzymujemy więc zamknięty szlak. Po jego usunięciu
równość stopni wejściowych i wyjściowych pozostaje prawdziwa.
Jeżeli nie wykorzystaliśmy wszystkich łuków, to pewien
niewykorzystany łuk jest incydentny z wierzchołkiem
już występującym w naszym szlaku (ma w nim początek
lub koniec): po zapomnieniu kierunków stosujemy
argument o pierwszej krawędzi opuszczającej zbiór
wierzchołków szlaku, jak w dowodzie
twierdzenia~\ref{thm:tg-euler-circuit}.
W multigrafie pozostałych łuków przynajmniej jeden
z dwóch stopni tego wierzchołka jest dodatni.
Ponieważ stopnie wejściowy i wyjściowy są równe,
oba są dodatnie.
Z tego wierzchołka konstruujemy następny niepusty
zamknięty szlak i wplatamy go do poprzedniego.
Iteracja kończy się po wykorzystaniu wszystkich łuków.
\end{proof}

\begin{twierdzenie}[Otwarty skierowany szlak Eulera]
\label{thm:tg-euler-directed-trail}
Niech $D=(V,A,\partial_-,\partial_+)$ będzie skończonym
multigrafem skierowanym z co najmniej jednym łukiem
i niech $s,t\in V$ będą różnymi wierzchołkami.
Istnieje skierowany szlak Eulera od $s$ do $t$
wtedy i tylko wtedy, gdy zachodzą oba warunki:
\begin{enumerate}[label=(\roman*)]
\item wszystkie wierzchołki incydentne z łukami należą
do jednej słabo spójnej składowej;
\item $d^+(s)-d^-(s)=1$, $d^+(t)-d^-(t)=-1$ oraz
$d^+(v)=d^-(v)$ dla każdego $v\in V\setminus\{s,t\}$.
\end{enumerate}
\end{twierdzenie}

\begin{proof}
Jeśli istnieje taki szlak, to po zapomnieniu kierunków
łuków łączy on wszystkie wierzchołki incydentne z łukami,
więc zachodzi (i). Każde wystąpienie wierzchołka jako
punktu pośredniego szlaku odpowiada jednemu wejściu
i jednemu wyjściu. Początkowe wyjście z $s$ i końcowe
wejście do $t$ są jedynymi niesparowanymi wystąpieniami.
Ponieważ szlak wykorzystuje każdy łuk dokładnie raz,
otrzymujemy wszystkie równości z (ii).

Załóżmy (i) i (ii). Dołączamy nowy łuk $a_*$ od $t$ do $s$,
różny od wszystkich dotychczasowych łuków, także wtedy,
gdy jeden z nich ma już początek $t$ i koniec $s$.
Nowy łuk zwiększa o jeden stopień wyjściowy $t$
oraz stopień wejściowy $s$, pozostawiając wszystkie
pozostałe stopnie bez zmian. W otrzymanym multigrafie
stopnie wejściowy i wyjściowy są więc równe przy każdym
wierzchołku. Wierzchołki $s$ i $t$ już należały do nośnika
krawędziowego na mocy (ii), więc (i) pozostaje spełniony.
Z twierdzenia~\ref{thm:tg-euler-directed-circuit}
otrzymujemy skierowany zamknięty szlak Eulera.

Zapis tego szlaku przesuwamy cyklicznie tak, aby jego
pierwszym wierzchołkiem było wystąpienie $s$ bezpośrednio
po przejściu przez $a_*$. Wówczas ostatni łuk zapisu
to $a_*$, prowadzący od $t$ do $s$. Po usunięciu tego
łuku i końcowego wystąpienia $s$ otrzymujemy skierowany
szlak od $s$ do $t$. Każdy pierwotny łuk występuje w nim
dokładnie raz, co kończy dowód.
\end{proof}

\begin{stwierdzenie}[Orientacja eulerowska]
Niech $G$ będzie skończonym multigrafem nieskierowanym.
Istnieje orientacja każdej jego krawędzi, dla której
$d^+(v)=d^-(v)$ przy wszystkich $v$, wtedy i tylko wtedy,
gdy każdy stopień $d_G(v)$ jest parzysty.
\end{stwierdzenie}

\begin{proof}
W każdej orientacji $d_G(v)=d^+(v)+d^-(v)$:
każda incydencja krawędzi z $v$ staje się początkiem
lub końcem łuku, przy czym pętla daje jedną incydencję
każdego rodzaju. Jeśli $d^+(v)=d^-(v)$, to
$d_G(v)=2d^+(v)$ jest parzyste.
Odwrotnie, w każdej nieizolowanej składowej grafu o
parzystych stopniach wybieramy cykl Eulera i orientujemy
krawędzie zgodnie z przebiegiem. Każde wejście na obiegu
odpowiada jednemu wyjściu. Wierzchołki izolowane mają
oba stopnie równe zero.
\end{proof}

\section{Silnie spójna orientacja i mosty}
\label{sec:tg-robbins}

\begin{twierdzenie}[Robbins]
\label{thm:tg-robbins}
Niech $G$ będzie skończonym spójnym grafem prostym
z co najmniej dwoma wierzchołkami. Istnieje silnie spójna
orientacja krawędzi $G$ wtedy i tylko wtedy,
gdy $G$ nie ma mostów.
\end{twierdzenie}

\begin{proof}
Jeżeli $e$ jest mostem, usunięcie go rozdziela graf na dwa
niepuste zbiory wierzchołków $S$ oraz $V\setminus S$.
Jedyny łuk łączący te części, niezależnie od wybranej
orientacji, umożliwia przejście tylko w jednym kierunku.
Silna spójność jest zatem niemożliwa.

Załóżmy, że $G$ nie ma mostów. Wykonujemy przeszukiwanie
w głąb (DFS) zaczynając od wierzchołka $r$ i otrzymujemy
drzewo rozpinające $T$. Orientujemy każdą krawędź drzewa
\emph{od korzenia ku potomkom}. Każdą pozostałą krawędź
orientujemy od potomka do przodka. Takie ustalenie
jest możliwe, ponieważ w nieskierowanym DFS każda
krawędź niedrzewowa łączy wierzchołek z jego przodkiem.
Wykonujemy DFS rekurencyjnie: po odkryciu wierzchołka
przeglądamy wszystkich jego sąsiadów, a dla każdego
dotąd nieodkrytego sąsiada zapisujemy krawędź odkrycia
i natychmiast wywołujemy DFS, oznaczając sąsiada jako
odkrytego na początku tego wywołania. Każdy wierzchołek
jest odkrywany najwyżej raz, więc procedura się kończy.
Gdyby po jej zakończeniu pozostał nieodkryty wierzchołek,
droga od $r$ do niego zawierałaby krawędź od odkrytego
do nieodkrytego wierzchołka; podczas przeglądania
sąsiadów pierwszego z nich odkrylibyśmy drugi.
Wszystkie wierzchołki zostają więc odkryte.
Każda zapisana krawędź łączy już odkryty wierzchołek
z nowym, toteż dołączenie jej zachowuje spójność
i nie tworzy cyklu. Zapisane krawędzie rzeczywiście
tworzą drzewo rozpinające $T$.
Aby uzasadnić własność krawędzi niedrzewowych, rozważmy
krawędź $xy$, przy czym $x$ odkryto przed $y$.
Wywołanie dla $x$ nie może zakończyć się przed odkryciem
$y$: podczas tego wywołania sąsiad $y$ zostałby wówczas
napotkany jako nieodkryty i odkryty rekurencyjnie.
Zatem $y$ odkryto podczas aktywnego wywołania dla $x$,
a wszystkie odkrycia w tym czasie pochodzą z wywołań
zagnieżdżonych w wywołaniu dla $x$. Stąd $y$ jest
potomkiem $x$ w drzewie DFS.

Z korzenia $r$ można dotrzeć do każdego wierzchołka
po skierowanych krawędziach drzewa. Pokażemy, że
z każdego wierzchołka można również dotrzeć do $r$.
Dla wierzchołka $v\ne r$ niech $p$ będzie jego rodzicem,
a $S_v$ zbiorem wszystkich potomków $v$ łącznie z $v$.
Krawędź $pv$ jest jedyną krawędzią drzewa łączącą
$S_v$ z $V\setminus S_v$: rodzic każdego innego
wierzchołka $S_v$ także należy do $S_v$, a każde
dziecko wierzchołka $S_v$ jest jego potomkiem.
Ponieważ $pv$ nie jest mostem, w $G-pv$ nadal istnieje
droga od $v$ do $p$. Musi ona zawierać krawędź między
$S_v$ i $V\setminus S_v$. Jest to krawędź niedrzewowa,
która na mocy własności DFS łączy pewnego potomka
$x\in S_v$
z przodkiem $y$ wierzchołka $v$ (ewentualnie $y=p$).
W orientacji możemy przejść od $v$ do $x$ w dół drzewa,
a potem łukiem $x\to y$. W ten sposób dotarliśmy do
wierzchołka o ściśle mniejszej głębokości niż $v$.
Powtarzając argument, osiągamy $r$, gdyż głębokość
jest nieujemną liczbą całkowitą i za każdym razem maleje.
Stąd dla dowolnych $u,w$ istnieje skierowany spacer
od $u$ przez $r$ do $w$. Jeżeli pewien wierzchołek
występuje w nim dwukrotnie, usuwamy fragment między
tymi wystąpieniami: pozostawione łuki nadal mają
właściwe kierunki i łączą te same końce spaceru.
Powtarzając tę operację, otrzymujemy skierowaną drogę. Orientacja jest silnie spójna.
\end{proof}

\begin{przyklad}[Dlaczego obecność mostu przeszkadza]
Dwa trójkąty połączone pojedynczą krawędzią nie mają
silnie spójnej orientacji: każda podróż między trójkątami
musi użyć mostu, którego orientacja dopuszcza tylko jeden
kierunek. Z drugiej strony dwa trójkąty o wspólnym
wierzchołku nie mają mostów i można je zorientować
zgodnie z obiegiem każdego trójkąta, otrzymując
silnie spójny graf skierowany.
\end{przyklad}

\begin{center}
\begin{tikzpicture}[scale=.86,line cap=round]
\foreach \x/\y/\n in {-3/0/a,-2/1/b,-1/0/c,1/0/d,2/1/e,3/0/f}
  \node[tgBlueVertex] (\n) at (\x,\y) {$\n$};
\draw[tgTree] (a)--(b)--(c)--(a) (d)--(e)--(f)--(d);
\draw[tgAccent,line width=2.8pt] (c)--(d);
\node[font=\footnotesize,text=tgCoral] at (0,.43) {most};
\node[font=\footnotesize,text=tgNavy] at (0,-.73)
 {Jedna krawędź między blokami uniemożliwia silną spójność orientacji.};
\end{tikzpicture}
\end{center}

